\documentclass{tietreport}

% Import images
\usepackage{graphicx}

% Bibliography configuration
\usepackage[
  date=year,
  style=alphabetic,
  backend=biber,
  sorting=nty,
  maxnames=5,
  minnames=3,
  maxalphanames=4,
  minalphanames=3,
  backref=true,
  doi=false,
  isbn=false,
  url=false,
  eprint=false
]{biblatex}
\addbibresource{references.bib}

%% ==================================================
%% CUSTOMIZE THESE FIELDS FOR YOUR PROJECT
%% ==================================================

% Course/Lab header
\TitlePageHeader{%
  UCS503 - Software Engineering%
}

% Main project title
\ProjectTitle{%
  TiTalks%
}

% Subtitle or project description
\ProjectSubTitle{%
  A Verified, Campus-Centred Social Networking Platform for TIET Students%
}

% Student information (up to 4 students)
\author{%
  \texttt{1024030304} Bhavya Agarwal \\
  \texttt{1024030301} Pratham Ranka \\
  \texttt{102403030298} Hitesh Arora \\
  \texttt{1024030299} Ananya Yadav%
}

% Group number, degree, year, program
\TitlePageSubText{%
  BE Third Year, Computer Science and Engineering%
}

% Faculty advisor/guide name
\AdvisorName{%
  Ms. Anuskha%
}

% Evaluation period and department info
\TitlePageFooterText{{%
    \large Project Progress Report \\[0.2cm]
    \bfseries Computer Science and Engineering
    Department%
  } \\[0.15cm]
  Thapar Institute of Engineering and Technology,
  Patiala%
}

%% ==================================================

\begin{document}

% Generate title page
\maketitle

% Table of contents (auto-generated from chapters)
\tableofcontents

% ===== CHAPTER 1: INTRODUCTION =====

\chapter{Introduction}

\section{Background and Context}

TiTalks is a campus-centred social networking application designed for the
students of Thapar Institute of Engineering and Technology (TIET). General
social media platforms contain a large amount of unrelated content and do not
provide a focused space for campus announcements, student activities,
societies, collaboration, and peer interaction. TiTalks addresses this gap by
providing a familiar visual social experience restricted conceptually to the
TIET community.

The current project is a feature-rich client-side prototype. It demonstrates
the complete intended user experience through responsive React screens,
typed mock data, local state, browser storage, local media, and direct URL
routing. The implemented prototype supports sign-in and onboarding, feeds,
stories, posts, reels, search, event discovery, messages, notifications,
profiles, settings, content management, and responsive navigation. A real
server, database, OAuth authentication, permanent media storage, and realtime
communication are intentionally reserved for the next development phase.

\subsection{Problem Statement}

TIET students currently depend on several unrelated channels to discover
campus events, follow societies, find peers, exchange updates, and maintain
professional or social connections. Information is fragmented, important
announcements compete with general-purpose content, and identity is not
necessarily tied to the institute. The project therefore aims to design and
prototype a single, accessible, and responsive campus network in which users
can discover relevant people and activities through an institute-focused
interface.

The immediate engineering problem is to validate the product flow before
investing in backend infrastructure. The prototype must behave like a
connected application during a demonstration while clearly separating mock
client-side behavior from features that require secure server-side support.

\section{Project Objectives}

\begin{enumerate}
  \item \textbf{Build the complete frontend experience:} Implement responsive
    desktop and mobile interfaces for authentication, feed browsing, content
    creation, discovery, messaging, notifications, profiles, reels, and
    settings.
  \item \textbf{Restrict the prototype to the TIET context:} Validate
    institutional email addresses under the \texttt{thapar.edu} domain and use
    campus-relevant societies, events, locations, and sample content.
  \item \textbf{Support realistic interaction flows:} Allow users to like,
    save, comment, share, follow, create, edit, archive, restore, delete,
    report, block, and mute content or accounts through local state.
  \item \textbf{Provide a maintainable architecture:} Use reusable typed
    components, route-level code splitting, feature-oriented folders, and a
    clear frontend-to-backend handoff plan.
  \item \textbf{Verify software quality:} Maintain a production build and an
    automated test pipeline containing unit, component, integration, browser
    smoke, responsive, and staging-style checks.
\end{enumerate}

\section{Scope and Current Status}

The present deliverable is a working frontend prototype rather than a fully
connected production system. Table~\ref{tab:status} distinguishes implemented
work from future backend work.

\begin{table}[h]
\centering
\begin{tabular}{|p{0.30\textwidth}|p{0.20\textwidth}|p{0.38\textwidth}|}
  \hline
  \textbf{Area} & \textbf{Status} & \textbf{Current Evidence} \\
  \hline
  Responsive frontend and navigation & Implemented & Desktop sidebar,
  collapsible right rail, mobile bottom navigation, and routed views \\
  \hline
  Social interactions & Implemented locally & Posts, comments, likes, saves,
  follows, stories, reels, messages, notifications, and content controls \\
  \hline
  Identity and preferences & Prototype complete & Institutional-email checks,
  onboarding, local session, profile settings, and persistent theme \\
  \hline
  Backend API and database & Planned & Node.js, Express, and PostgreSQL are
  documented for the connected MVP \\
  \hline
  Production authentication and uploads & Planned & Google OAuth, server
  sessions, authorization, and durable media storage are not yet connected \\
  \hline
\end{tabular}
\caption{Current TiTalks implementation status}
\label{tab:status}
\end{table}

% ===== CHAPTER 2: METHODOLOGY & DESIGN =====

\chapter{Methodology and Design}

\section{System Architecture}

The project follows an incremental prototype-first methodology. The team first
defined user flows and built a complete client experience using local data.
This approach makes navigation, interaction design, responsiveness, and
feature scope testable before introducing deployment and data-management
complexity. The current architecture is a single-page application executed in
the browser.

\subsection{High-Level Design}

The application is divided into four logical layers:

\begin{enumerate}
  \item \textbf{Presentation layer:} React views and reusable UI components
    render the feed, profiles, discovery pages, inbox, activity, reels, and
    modal workflows.
  \item \textbf{Navigation layer:} React Router maps direct browser URLs to
    screens and supports browser history and deployment rewrites.
  \item \textbf{Client state layer:} React state stores temporary posts,
    stories, reels, reactions, follows, archive state, and interface feedback.
    Browser local storage preserves the mock session, profile settings, and
    selected appearance theme.
  \item \textbf{Data and media layer:} Typed mock records and bundled WebP/MP4
    assets provide deterministic demonstration data without requiring a
    backend connection.
\end{enumerate}

The planned connected architecture will retain the React frontend and replace
mock state with HTTP/JSON calls to a Node.js and Express application backed by
PostgreSQL. Business rules, authorization, persistence, and media ownership
will move to the server at that stage.

\subsection{Technical Stack}

\begin{itemize}
  \item \textbf{Framework and language:} React 18 with TypeScript provides a
    component model and compile-time checking for shared data contracts
    \cite{react,typescript}.
  \item \textbf{Build tool:} Vite provides the development server, production
    bundling, and route-level chunk generation \cite{vite}.
  \item \textbf{Styling:} Tailwind CSS v4 is used through component-local
    utility classes, with a persistent light/dark design system
    \cite{tailwind}.
  \item \textbf{Routing and icons:} React Router manages application URLs and
    Lucide React supplies consistent interface icons.
  \item \textbf{Motion and scrolling:} Browser animation APIs and Lenis are
    used with reduced-motion safeguards.
  \item \textbf{Testing:} Vitest, Testing Library, jsdom, Playwright Core, and
    custom Node.js deployment checks form the automated quality pipeline
    \cite{playwright}.
  \item \textbf{External data:} The weather widget requests current Patiala
    conditions from Open-Meteo and includes loading and error states
    \cite{openmeteo}.
  \item \textbf{Planned backend:} Node.js, Express, PostgreSQL, Google OAuth,
    and durable media storage.
\end{itemize}

\section{Implementation Details}

\subsection{Module 1: Authentication and User Onboarding}

The authentication interface supports mock Google account selection, direct
sign-in, registration, password recovery navigation, and profile onboarding.
The reusable \texttt{isTietEmail} utility normalizes an address and accepts the
\texttt{thapar.edu} domain and its subdomains while rejecting personal or
misleading domains. Password length and required profile fields are validated
before the local authentication callback is executed. The current session is
stored under a versioned local-storage key so that the browser demo remains
usable after refresh.

This module demonstrates the intended product flow but does not claim to be a
secure authentication system. Password verification, Google OAuth tokens,
server sessions, rate limiting, and authorization checks require the planned
backend.

\subsection{Module 2: Feed, Posts, Stories, and Reels}

The central feed presents campus updates with author information, location,
caption, tags, engagement counts, and media. Users can like, save, share,
comment, open routed details, and access context-sensitive content controls.
The composer supports local image selection and immediate insertion of posts,
stories, or reels into client state. Story navigation supports keyboard input,
while story management includes creation, replacement, deletion, and highlight
selection. The reels view provides real local video playback, interaction
controls, creation, editing, deletion, and incremental feed loading.

\subsection{Module 3: Discovery, Profiles, and Campus Context}

The People and Pulse views support profile discovery and event exploration.
Search filters local typed records, follow actions update local relationship
state, and other-user profile routes display dedicated profile experiences.
The Pulse area supports category filtering, event details, adding an event to
the user's day, and a responsive TIET campus map viewer. The home page also
contains a live Patiala weather component with explicit loading, success, and
failure behavior.

The profile module contains post, saved, and tagged grids; follower and
following lists; highlights; archived posts; profile editing; privacy and
notification controls; password recovery; logout; and account-deletion UI.

\subsection{Module 4: Messaging, Notifications, and Safety Flows}

The inbox supports direct and group conversation creation, text messages,
image messages, selection state, and automatic scrolling to recent messages.
The activity view represents notifications and follow requests. Content menus
provide edit, archive, restore, delete, report, block, mute, and unfollow
flows. These operations currently modify local interface state; enforceable
privacy and moderation rules will require server-side identity and ownership
checks.

\subsection{Module 5: Responsive Design and Accessibility}

The interface adapts between desktop and mobile layouts. Desktop users receive
an expandable sidebar and collapsible agenda rail, while mobile users receive
bottom navigation and touch-friendly controls. Loading, empty, error, and
not-found states are represented explicitly. Buttons use accessible labels,
status feedback uses live regions, keyboard dismissal is provided for relevant
dialogs, and animations respect the user's reduced-motion preference in line
with accessibility guidance \cite{wcag}.

\subsection{Development Process}

Work was organized incrementally by feature. Reusable primitives and shared
types were established first, followed by core feed and navigation behavior,
then authentication, profiles, discovery, messaging, safety controls, reels,
responsive refinements, and automated regression checks. Route-level views are
loaded lazily, which keeps optional screen code outside the initial application
bundle until required.

% ===== CHAPTER 3: RESULTS & EVALUATION =====

\chapter{Results and Evaluation}

\section{Testing Methodology}

The project now uses a layered automated testing strategy. Each layer checks a
different class of failure, and the complete pipeline is available through
\texttt{npm run test:all}.

\begin{itemize}
  \item \textbf{Unit testing:} Vitest verifies the institutional-email utility
    against valid, mixed-case, subdomain, personal-domain, malformed, and empty
    inputs. Component-level unit tests verify the theme toggle and reusable
    loading/error state panel.
  \item \textbf{Integration testing:} Testing Library renders the
    authentication screen inside a memory router. Tests verify rejected
    personal-email sign-in, successful normalized institutional sign-in, and
    the complete registration-to-profile-onboarding flow.
  \item \textbf{Build and static verification:} TypeScript strict checking is
    executed without emitting files, followed by a Vite production build. This
    catches type incompatibilities, missing imports, and bundling failures.
  \item \textbf{Browser smoke and system testing:} A headless Chromium test
    exercises a realistic production build. It verifies email rejection,
    authentication state, home content, weather rendering, sidebar geometry,
    right-rail controls, campus map behavior, Pulse filters, event planning,
    Reels loading and interaction, creation modal access, and mobile rendering.
  \item \textbf{Staging-style testing:} A temporary static server hosts the
    production bundle. Automated checks confirm the root page, generated
    JavaScript and CSS assets, and SPA fallback behavior for seven deep routes:
    People, Pulse, Reels, Inbox, Activity, Profile, and Settings.
  \item \textbf{Regression gate:} The combined command runs all automated
    tests, the production build, browser smoke checks, and staging checks in a
    fixed order. A failure at any stage prevents the pipeline from reporting
    success.
\end{itemize}

\section{Performance Metrics}

The latest verified execution completed successfully. The pass percentage
below is the test pass rate for the currently defined automated cases; it is
not a claim of 100\% source-code coverage.

\begin{table}[h]
\centering
\begin{tabular}{|p{0.30\textwidth}|p{0.22\textwidth}|p{0.34\textwidth}|}
  \hline
  \textbf{Metric} & \textbf{Target} & \textbf{Achieved} \\
  \hline
  Automated test cases & All pass & 15 of 15 passed (100\% pass rate) \\
  \hline
  Test files & All pass & 4 of 4 passed \\
  \hline
  TypeScript check & Zero errors & Passed with zero type errors \\
  \hline
  Production build & Successful bundle & Passed; 1,628 modules transformed \\
  \hline
  Browser smoke test & Critical flows pass & Passed on desktop and mobile
  viewports \\
  \hline
  Staging deep routes & Seven routes available & 7 of 7 routes passed \\
  \hline
  Built entry assets & JavaScript and CSS available & 2 of 2 referenced entry
  assets passed \\
  \hline
  Browser page errors & Zero during smoke flow & Zero detected \\
  \hline
\end{tabular}
\caption{Verified quality-gate results}
\label{tab:performance}
\end{table}

The production output uses route-level chunks for major screens. The generated
entry JavaScript bundle is approximately 301.54~kB before compression and
91.25~kB after gzip compression. The generated stylesheet is approximately
71.52~kB before compression and 11.59~kB after gzip compression. These figures
provide a reproducible baseline for later optimization; they are build sizes,
not network latency measurements.

\section{Analysis}

The current implementation meets its primary frontend-prototype objectives.
The application is demonstrable without a server, offers broad feature
coverage, and behaves consistently across desktop and mobile layouts. Typed
models and reusable component groups provide a practical foundation for
replacing mock data with APIs. Automated testing now covers isolated logic,
component behavior, cross-component authentication flows, production-browser
behavior, and deployment routing.

The strongest result is the completeness of the user experience at the
prototype stage. A presenter can move from authentication to feed interaction,
discovery, content creation, messaging, profile management, reels, safety
actions, and responsive layouts without unfinished screens. Local media also
makes the main demo independent of network availability, except for the live
weather enhancement, which has a visible error state and is mocked during
browser testing.

The principal limitation is that client-side behavior cannot provide durable
or secure multi-user functionality. Refreshing the page resets many social
interactions, mock identities do not prove real ownership, and local report or
block actions cannot be enforced across accounts. Performance under concurrent
users, database correctness, API authorization, and realtime messaging have
not yet been measured because those layers do not exist in the current phase.

The testing result is therefore interpreted carefully: all implemented
automated tests pass, but additional backend, security, load, and end-to-end
tests must be added when the connected system is developed.

% ===== CHAPTER 4: CONCLUSION =====

\chapter{Conclusion}

\section{Summary of Work}

TiTalks has progressed from a campus-network concept to a complete interactive
frontend prototype. The team implemented institutional-email validation,
authentication and onboarding screens, a responsive feed, stories, reels,
content creation and management, discovery, profiles, messaging,
notifications, campus event features, settings, theming, motion, and robust
empty/error states. The codebase is organized into typed components, views,
modals, layouts, data models, and utilities.

The project also includes a repeatable quality gate. At the time of this
report, all 15 automated tests passed, the TypeScript production build
succeeded, the browser smoke flow passed, and the staging-style route and asset
checks passed. This establishes a reliable baseline for the backend phase.

\section{Key Findings}

\begin{itemize}
  \item \textbf{Finding 1:} A prototype-first process allowed the team to
    validate a large number of campus social workflows before committing to a
    server and database design.
  \item \textbf{Finding 2:} Shared TypeScript models, feature-based components,
    and route-level lazy loading support maintainability despite the breadth of
    the interface.
  \item \textbf{Finding 3:} Layered tests reveal different defects: unit tests
    protect validation rules, integration tests protect user flows, the build
    protects type and bundle correctness, browser smoke tests protect real UI
    behavior, and staging checks protect deployment routing.
  \item \textbf{Finding 4:} The frontend is ready for API integration, but it
    must not be described as production-ready until identity, authorization,
    persistence, moderation, and upload security are implemented server-side.
\end{itemize}

\section{Future Work and Recommendations}

\begin{enumerate}
  \item Implement Google OAuth restricted to TIET accounts and maintain secure
    server-side sessions.
  \item Build a Node.js and Express REST API for users, profiles, posts,
    comments, reactions, saves, follows, search, and discovery.
  \item Add PostgreSQL persistence with migrations, constraints, indexes, and
    transactional tests.
  \item Introduce validated media uploads, object storage, ownership checks,
    file-size limits, and content-type restrictions.
  \item Add realtime direct and group messaging, notification delivery, and
    read-state synchronization.
  \item Enforce privacy, block, mute, report, archive, and deletion rules on the
    server rather than only in interface state.
  \item Extend the automated suite with API unit tests, database integration
    tests, authorization tests, security checks, accessibility scans, load
    tests, and full user-acceptance scenarios.
  \item Add continuous integration so every proposed change runs
    \texttt{npm run test:all} before merging or deployment.
  \item Measure real-user performance and accessibility after deployment, then
    optimize the largest route chunks and media delivery where evidence shows
    a need.
\end{enumerate}

\section{Final Remarks}

The present TiTalks prototype successfully demonstrates the intended campus
experience and provides enough interaction depth for evaluation and user
feedback. The next milestone is not to redesign the interface, but to connect
the validated flows to secure, persistent backend services. The project's
current separation between UI behavior and planned server responsibilities
gives the team a clear and manageable path toward a deployable campus platform.

% ===== BIBLIOGRAPHY =====

% Print bibliography with heading in TOC
\printbibliography[heading=bibintoc]

\end{document}
